\chapter{Konzeption und prototypische Implementierung eines VR-basierten Tools zur kollaborativen Annotation von BIM-Modellen auf Basis der Unreal Engine mit IFC-/BCF-basierter Persistierung und Integration in OpenBIM-Workflows}

\section{Einleitung und Motivation}

Die Digitalisierung der Bauindustrie hat in den letzten Jahren durch den Einsatz von Building Information Modeling (BIM) erheblich an Bedeutung gewonnen.
BIM beschreibt die digitale Repräsentation eines Bauwerks unter Integration geometrischer, semantischer und funktionaler Informationen und dient als zentrale Grundlage für Planung, Ausführung und Betrieb von Bauprojekten.\cite{Eastman2011}
Parallel dazu hat sich Virtual Reality (VR) als Schlüsseltechnologie zur immersiven Visualisierung komplexer Daten etabliert.
Insbesondere im Bauwesen ermöglicht VR eine realitätsnahe Begehung von Gebäudemodellen bereits in frühen Planungsphasen, wodurch Planungsfehler frühzeitig erkannt und Kommunikationsprozesse verbessert werden können.
Planer, Auftraggeber und Öffentlichkeit können Projekte realistischer beurteilen, was zu fundierteren Entscheidungen führt und Fehlinterpretationen reduziert.
Aktuelle Forschungsarbeiten zeigen, dass VR-basierte Systeme im Bauwesen inzwischen weit über reine Demonstratoren hinausgehen und insbesondere für Design Reviews, Planungsabstimmungen und die Kommunikation zwischen Stakeholdern eingesetzt werden. Gleichzeitig wird deutlich, dass viele Lösungen weiterhin primär auf die Visualisierung von BIM-Modellen ausgerichtet sind, während kollaborative Interaktion, Annotation und die strukturierte Rückführung von Erkenntnissen in BIM-Prozesse noch nicht durchgängig unterstützt werden.
VR birgt also großes Potenzial, Planungsprozesse zu verbessern und transparenter zu machen und als leistungsfähiges Werkzeug verständliche, partizipative und realistische Planung zu ermöglichen.\cite{Portman2015}

Trotz dieser Entwicklungen zeigt sich ein grundlegendes Defizit in der praktischen Anwendung: Während BIM-Systeme detaillierte Modellierung und Datenverwaltung ermöglichen, werden VR-Anwendungen meist lediglich zur Visualisierung eingesetzt.
Eine direkte Interaktion mit dem Modell – insbesondere im Sinne kollaborativer Bearbeitung oder strukturierter Annotation – ist nur eingeschränkt möglich.\cite{Whyte2018}
Dies führt zu einem fragmentierten Workflow: Erkenntnisse aus VR-Begehungen müssen manuell dokumentiert und anschließend in das BIM-System übertragen werden.
Dieser Medienbruch verursacht Ineffizienz, erhöht die Fehleranfälligkeit und erschwert die Zusammenarbeit zwischen Projektbeteiligten wie Architekten, Ingenieuren und anderen Stakeholdern.
Ein besonders praxisrelevantes Szenario ist die Präsentation eines Bauprojekts beim Kunden vor Ort.
Hierbei könnte eine VR-Brille genutzt werden, um das Modell immersiv zu erleben und Änderungswünsche direkt zu erfassen.
Obwohl mittlerweile zahlreiche kommerzielle und wissenschaftliche Lösungen die Visualisierung von BIM-Modellen in VR unterstützen, bestehen weiterhin Defizite bei der kollaborativen Erfassung von Feedback, der strukturierten Verwaltung von Annotationen sowie deren Integration in bestehende OpenBIM-Workflows.
Insbesondere die direkte Verknüpfung von VR-basierten Anmerkungen mit IFC-basierten Modellinformationen stellt weiterhin eine Herausforderung dar.

\section{Stand der Forschung und Problemstellung}

Die Kombination von BIM und VR wurde in der Forschung bereits intensiv untersucht.
Studien zeigen, dass VR insbesondere das räumliche Verständnis verbessert und die Kommunikation zwischen Projektbeteiligten unterstützt.\cite{Whyte2018}
Du et al.\ zeigen, dass VR-basierte Systeme die Effizienz von Design-Reviews erhöhen können.
Gleichzeitig wird jedoch deutlich, dass die meisten Systeme keine bidirektionale Integration mit BIM-Daten ermöglichen.
Kollaborative VR-Umgebungen wurden unter anderem im Kontext von Multi-User-Systemen untersucht.
Messner et al.\ beschreiben frühe Ansätze, bei denen mehrere Nutzer gleichzeitig auf ein virtuelles Gebäudemodell zugreifen können.
Neuere Arbeiten erweitern diesen Ansatz in Richtung „Metaverse“-ähnlicher Systeme, bei denen persistente virtuelle Räume entstehen.\cite{Wang2018}
Dennoch bleibt die Integration in bestehende BIM-Workflows eine offene Herausforderung.
Es fehlt ein integriertes System, das eine kollaborative Interaktion mit BIM-Modellen in einer VR-Umgebung ermöglicht und dabei eine strukturierte Rückführung von Annotationen in den BIM-Workflow unterstützt.
Annotationen sind ein zentrales Element in kollaborativen Planungsprozessen.
In klassischen BIM-Workflows werden hierfür häufig externe Tools oder Formate wie das BIM Collaboration Format (BCF) verwendet.
In VR-Umgebungen existieren zwar Ansätze zur Platzierung von Markern oder Kommentaren, jedoch fehlt in der Regel eine strukturierte Verknüpfung mit BIM-Elementen sowie eine automatisierte Rückführung in das Ursprungsmodell.\cite{Wang2018}
Bestehende Lösungen weisen mehrere Einschränkungen auf:

\begin{itemize}
\item VR-Anwendungen sind primär visualisierungsorientiert und erlauben keine tiefgehende Interaktion mit BIM-Daten.\cite{Du2018}
\item BIM-Modelle werden häufig in polygonale Darstellungen überführt, wodurch semantische Informationen verloren gehen.\cite{Wang2018}
\item Kollaborative Funktionen sind entweder nicht vorhanden oder nicht mit BIM-Systemen integriert.\cite{Messner2003}
\end{itemize}

Darüber hinaus besteht eine Diskrepanz zwischen Forschung und Praxis: Während prototypische Ansätze zur Integration von VR und BIM existieren, fehlen praxistaugliche, leicht einsetzbare Lösungen für den mobilen Einsatz beim Kunden.

\section{Zielsetzung der Arbeit}

Ziel dieser Bachelorarbeit ist die Konzeption und prototypische Implementierung eines Systems, das die genannten Defizite adressiert.
Der Fokus liegt dabei bewusst auf der Annotation und Interaktion, nicht auf der vollständigen Bearbeitung von BIM-Modellen in VR.\@
Im Zuge der Arbeit wird zunächst ein Entwurf für ein VR-gestütztes System zur Darstellung von BIM-Modellen erstellt.
Auf dieser Basis wird ein funktionierender Prototyp unter Einsatz der Unreal Engine als technische Grundlage entwickelt.
Zentraler Schwerpunkt der Arbeit ist die Entwicklung eines kollaborativen Mehrbenutzersystems, das mehreren Anwendern die gleichzeitige Betrachtung eines BIM-Modells innerhalb einer VR-Umgebung ermöglicht.
Nutzer sollen Annotationen direkt an Modellelementen platzieren, gegenseitig einsehen und gemeinsam bearbeiten können.
Die erzeugten Annotationen sollen dauerhaft gespeichert und über IFC-Elementreferenzen beziehungsweise das BIM Collaboration Format (BCF) in bestehende BIM-Prozesse integriert werden können.
Zusätzlich werden unterschiedliche Methoden zur systematischen Speicherung und zur anschließenden Verarbeitung der gesammelten Annotationen analysiert und beurteilt.
Als Ergebnis der Arbeit wird ein funktionsfähiger VR-Prototyp erwartet.
Darüber hinaus soll ein Konzept entwickelt und prototypisch umgesetzt werden, das Annotationen eindeutig mit IFC-Objekten verknüpft und diese Informationen für nachgelagerte BIM-Prozesse bereitstellt.
Hierzu werden insbesondere BCF-basierte Ansätze zur Speicherung und zum Austausch von Issues untersucht und bewertet.
Bestehende BIM-Workflows werden dabei als Prozesskette aus IFC-Modellierung, VR-basierter Begutachtung, kollaborativer Annotation sowie anschließender Weiterverarbeitung in BIM-Autoren- oder Koordinationssystemen verstanden.
Ergänzend dazu erfolgt eine Analyse der technischen Herausforderungen, die im Rahmen der Umsetzung und Integration auftreten.

\section{Vorläufige Gliederung}

\begin{enumerate}
\item Einleitung
\item Grundlagen
\item Stand der Forschung
\item Anforderungsanalyse und Systemkonzept
\item Entwurf und Implementierung
\item Evaluation und Diskussion
\item Fazit und Ausblick
\end{enumerate}